% GENERATED FILE. Edit canonical lessons, then run npm run generate:books.
% canonical-source-hash: fnv1a64:0b32c8bcf1a3d6b2
% canonical-lessons: HI-C559-kalpna, HI-C559-suruat, HI-C559-udghatan, HI-C559-bhavi, HI-C559-zindagi, HI-R559-review-553-559, HI-R559-review-final-559

\chapter{Looking Ahead}
\label{ch:hi-looking-ahead}
\hlchaptermodality{559}

\begin{chapteropening}
\textbf{By the end of this chapter:} \emph{I can talk about beginnings, what I imagine and the life ahead.}

Complete the last lesson of chapter 559: I can talk about beginnings, what I imagine and the life ahead.
\end{chapteropening}

\section[\texorpdfstring{\hi{कल्पना} (kalpnā)}{कल्पना (kalpnā)}]{\hi{कल्पना} (kalpnā) — imagination; an idea of what could be}

\label{lesson:HI-C559-kalpna}



\begin{quote}
Before the new one: say the Hindi for an ambition, then the Hindi for a passion.
\end{quote}

\subsection*{You'll want to know: \hi{कल्पना}}
\textbf{\hi{कल्पना}} — \emph{kalpnā} — \textquotedblleft{}imagination; an idea of what could be\textquotedblright{}.

\begin{quote}
\hi{बच्चों} \hi{की} \hi{कल्पना} \hi{बहुत} \hi{तेज़} \hi{होती} \hi{है।}
\end{quote}

\begin{quote}
\emph{Baccoṁ kī kalpnā bahut tez hotī hai.}
\end{quote}

\begin{quote}
Children have very lively imaginations.
\end{quote}

Feminine.

\begin{cousinweb}
\textbf{\hi{कल्पना}} (\emph{kalpnā}) is a tatsama, taken whole from Sanskrit \emph{kalpanā}, \textquotedblleft{}making; creating in the mind, fiction\textquotedblright{} (Monier-Williams).
\end{cousinweb}

\begin{grammarlens}[title={What you've built}]
One more word for saying what happens if.
\end{grammarlens}

\subsection*{Your turn}
\begin{itemize}
\raggedright
  \item \emph{Say it:} \emph{kalpnā}
  \item \emph{Say it:} \emph{kalpnā}, once more
  \item \emph{Say it:} \emph{Baccoṁ kī kalpnā bahut tez hotī hai.} — Children have very lively imaginations.
  \item \emph{Recall:} say the Hindi for an ambition, then the Hindi for a passion, then say \emph{kalpnā} again
\end{itemize}

\begin{tcolorbox}[breakable,colback=teal!4,colframe=teal!35!black,title={Before you move on}]
What does \hi{कल्पना} mean? (\textquotedblleft{}Imagination; an idea of what could be\textquotedblright{}.) Say it once more.
\end{tcolorbox}
\section[\texorpdfstring{\hi{शुरुआत} (śuruāt)}{शुरुआत (śuruāt)}]{\hi{शुरुआत} (śuruāt) — a beginning; a start}

\label{lesson:HI-C559-suruat}



\begin{quote}
Before the new one: say the Hindi for a passion, then the Hindi for imagination.
\end{quote}

\subsection*{You'll want to know: \hi{शुरुआत}}
\textbf{\hi{शुरुआत}} — \emph{śuruāt} — \textquotedblleft{}a beginning; a start\textquotedblright{}.

\begin{quote}
\hi{नए} \hi{साल} \hi{की} \hi{शुरुआत} \hi{अच्छी} \hi{हुई।}
\end{quote}

\begin{quote}
\emph{Nae sāl kī śuruāt acchī huī.}
\end{quote}

\begin{quote}
The new year got off to a good start.
\end{quote}

Feminine.

\begin{cousinweb}
\textbf{\hi{शुरुआत}} (\emph{śuruāt}) is usually traced to Arabic \emph{shurūʿ}, \textquotedblleft{}a beginning\textquotedblright{}, through Persian. Hindi adds -āt to make a noun of it.
\end{cousinweb}

\begin{grammarlens}[title={What you've built}]
One more word for saying what happens if.
\end{grammarlens}

\subsection*{Your turn}
\begin{itemize}
\raggedright
  \item \emph{Say it:} \emph{śuruāt}
  \item \emph{Say it:} \emph{śuruāt}, once more
  \item \emph{Say it:} \emph{Nae sāl kī śuruāt acchī huī.} — The new year got off to a good start.
  \item \emph{Recall:} say the Hindi for a passion, then the Hindi for imagination, then say \emph{śuruāt} again
\end{itemize}

\begin{tcolorbox}[breakable,colback=teal!4,colframe=teal!35!black,title={Before you move on}]
What does \hi{शुरुआत} mean? (\textquotedblleft{}A beginning; a start\textquotedblright{}.) Say it once more.
\end{tcolorbox}
\section[\texorpdfstring{\hi{उद्घाटन} (udghāṭan)}{उद्घाटन (udghāṭan)}]{\hi{उद्घाटन} (udghāṭan) — an opening; an inauguration}

\label{lesson:HI-C559-udghatan}



\begin{quote}
Before the new one: say the Hindi for imagination, then the Hindi for a beginning.
\end{quote}

\subsection*{You'll want to know: \hi{उद्घाटन}}
\textbf{\hi{उद्घाटन}} — \emph{udghāṭan} — \textquotedblleft{}an opening; an inauguration\textquotedblright{}.

\begin{quote}
\hi{मंत्री} \hi{ने} \hi{नए} \hi{पुल} \hi{का} \hi{उद्घाटन} \hi{किया।}
\end{quote}

\begin{quote}
\emph{Mantrī ne nae pul kā udghāṭan kiyā.}
\end{quote}

\begin{quote}
The minister opened the new bridge.
\end{quote}

Masculine.

\begin{cousinweb}
\textbf{\hi{उद्घाटन}} (\emph{udghāṭan}) is a tatsama, taken whole from Sanskrit \emph{udghāṭana}, \textquotedblleft{}opening, unlocking\textquotedblright{} (Monier-Williams).
\end{cousinweb}

\begin{grammarlens}[title={What you've built}]
One more word for saying what happens if.
\end{grammarlens}

\subsection*{Your turn}
\begin{itemize}
\raggedright
  \item \emph{Say it:} \emph{udghāṭan}
  \item \emph{Say it:} \emph{udghāṭan}, once more
  \item \emph{Say it:} \emph{Mantrī ne nae pul kā udghāṭan kiyā.} — The minister opened the new bridge.
  \item \emph{Recall:} say the Hindi for imagination, then the Hindi for a beginning, then say \emph{udghāṭan} again
\end{itemize}

\begin{tcolorbox}[breakable,colback=teal!4,colframe=teal!35!black,title={Before you move on}]
What does \hi{उद्घाटन} mean? (\textquotedblleft{}An opening; an inauguration\textquotedblright{}.) Say it once more.
\end{tcolorbox}
\section[\texorpdfstring{\hi{भावी} (bhāvī)}{भावी (bhāvī)}]{\hi{भावी} (bhāvī) — future; to come}

\label{lesson:HI-C559-bhavi}



\begin{quote}
Before the new one: say the Hindi for a beginning, then the Hindi for an opening.
\end{quote}

\subsection*{You'll want to know: \hi{भावी}}
\textbf{\hi{भावी}} — \emph{bhāvī} — \textquotedblleft{}future; to come\textquotedblright{}.

\begin{quote}
\hi{भावी} \hi{डाक्टर} \hi{मेहनत} \hi{से} \hi{पढ़} \hi{रहे} \hi{हैं।}
\end{quote}

\begin{quote}
\emph{Bhāvī ḍākṭar mehnat se paṛh rahe haiṁ.}
\end{quote}

\begin{quote}
The doctors of the future are studying hard.
\end{quote}

\begin{cousinweb}
\textbf{\hi{भावी}} (\emph{bhāvī}) is a tatsama, taken whole from Sanskrit \emph{bhāvin}, \textquotedblleft{}becoming; future\textquotedblright{} (Monier-Williams).
\end{cousinweb}

\begin{grammarlens}[title={What you've built}]
One more word for saying what happens if.
\end{grammarlens}

\subsection*{Your turn}
\begin{itemize}
\raggedright
  \item \emph{Say it:} \emph{bhāvī}
  \item \emph{Say it:} \emph{bhāvī}, once more
  \item \emph{Say it:} \emph{Bhāvī ḍākṭar mehnat se paṛh rahe haiṁ.} — The doctors of the future are studying hard.
  \item \emph{Recall:} say the Hindi for a beginning, then the Hindi for an opening, then say \emph{bhāvī} again
\end{itemize}

\begin{tcolorbox}[breakable,colback=teal!4,colframe=teal!35!black,title={Before you move on}]
What does \hi{भावी} mean? (\textquotedblleft{}Future; to come\textquotedblright{}.) Say it once more.
\end{tcolorbox}
\section[\texorpdfstring{\hi{ज़िंदगी} (zindagī)}{ज़िंदगी (zindagī)}]{\hi{ज़िंदगी} (zindagī) — life (as one lives it)}

\label{lesson:HI-C559-zindagi}



\begin{quote}
Before the new one: say the Hindi for an opening, then the Hindi for future.
\end{quote}

\subsection*{You'll want to know: \hi{ज़िंदगी}}
\textbf{\hi{ज़िंदगी}} — \emph{zindagī} — \textquotedblleft{}life (as one lives it)\textquotedblright{}.

\begin{quote}
\hi{गाँव} \hi{की} \hi{ज़िंदगी} \hi{शहर} \hi{से} \hi{अलग} \hi{है।}
\end{quote}

\begin{quote}
\emph{Gā̃v kī zindagī śahar se alag hai.}
\end{quote}

\begin{quote}
Village life is different from city life.
\end{quote}

Feminine.

\begin{cousinweb}
\textbf{\hi{ज़िंदगी}} (\emph{zindagī}) is usually traced to Persian \emph{zindagī}, \textquotedblleft{}life\textquotedblright{}.
\end{cousinweb}

\begin{grammarlens}[title={What you've built}]
One more word for saying what happens if.
\end{grammarlens}

\subsection*{Your turn}
\begin{itemize}
\raggedright
  \item \emph{Say it:} \emph{zindagī}
  \item \emph{Say it:} \emph{zindagī}, once more
  \item \emph{Say it:} \emph{Gā̃v kī zindagī śahar se alag hai.} — Village life is different from city life.
  \item \emph{Recall:} say the Hindi for an opening, then the Hindi for future, then say \emph{zindagī} again
\end{itemize}

\begin{tcolorbox}[breakable,colback=teal!4,colframe=teal!35!black,title={Before you move on}]
What does \hi{ज़िंदगी} mean? (\textquotedblleft{}Life (as one lives it)\textquotedblright{}.) Say it once more.
\end{tcolorbox}
\section[(dialogue)]{Looking back: from In the Right Order to Looking Ahead}

\label{lesson:HI-R559-review-553-559}



\begin{quote}
Say the Hindi for except (\hi{के} \hi{सिवा}) and a series.
\end{quote}

\subsection*{Your turn}
Say these aloud:

\begin{itemize}
\raggedright
  \item except (\hi{के} \hi{सिवा}), a series, order, a stage, a phase, a date (formal)
  \item a lesson (learned), ultimately, as a result, approximately, actually
  \item to bear fruit, to sway, to rest (something) against, a fall, a downfall
  \item doubt (formal), certainly, surely, ever, a groundless fear, a firm decision
  \item a worry, peace of mind, ease, a guarantee, misfortune, fear, dread
  \item a wish, a longing (for something), crazy (about), an interest, an ambition, a passion
  \item imagination, a beginning, an opening, future, life (as one lives it)
\end{itemize}

\begin{tcolorbox}[breakable,colback=teal!4,colframe=teal!35!black,title={Before you move on}]
Except (\hi{के} \hi{सिवा})? (\textbf{\hi{सिवा}}, \emph{sivā}.) A series? (\textbf{\hi{सिलसिला}}, \emph{silsilā}.) Order? (\textbf{\hi{क्रम}}, \emph{kram}.) A stage, a phase? (\textbf{\hi{चरण}}, \emph{caraṇ}.) A date (formal)? (\textbf{\hi{तिथि}}, \emph{tithi}.) A lesson (learned)? (\textbf{\hi{सबक़}}, \emph{sabaq}.) Ultimately? (\textbf{\hi{अंततः}}, \emph{antataḥ}.) As a result? (\textbf{\hi{नतीजतन}}, \emph{natījatan}.) Approximately? (\textbf{\hi{तक़रीबन}}, \emph{taqrīban}.) Actually? (\textbf{\hi{दरअसल}}, \emph{darasal}.) To bear fruit? (\textbf{\hi{फलना}}, \emph{phalnā}.) To sway? (\textbf{\hi{डोलना}}, \emph{ḍolnā}.) To rest (something) against? (\textbf{\hi{टिकाना}}, \emph{ṭikānā}.) A fall? (\textbf{\hi{गिरावट}}, \emph{girāvaṭ}.) A downfall? (\textbf{\hi{पतन}}, \emph{patan}.) Doubt (formal)? (\textbf{\hi{संदेह}}, \emph{sandeh}.) Certainly, surely? (\textbf{\hi{यक़ीनन}}, \emph{yaqīnan}.) Ever? (\textbf{\hi{हरगिज़}}, \emph{hargiz}.) A groundless fear? (\textbf{\hi{वहम}}, \emph{vaham}.) A firm decision? (\textbf{\hi{निश्चय}}, \emph{niścay}.) A worry? (\textbf{\hi{अंदेशा}}, \emph{andeśā}.) Peace of mind, ease? (\textbf{\hi{इत्मीनान}}, \emph{itmīnān}.) A guarantee? (\textbf{\hi{गारंटी}}, \emph{gāraṇṭī}.) Misfortune? (\textbf{\hi{दुर्भाग्य}}, \emph{durbhāgya}.) Fear, dread? (\textbf{\hi{ख़ौफ़}}, \emph{khauf}.) A wish, a longing (for something)? (\textbf{\hi{ख़्वाहिश}}, \emph{khvāhiś}.) Crazy (about)? (\textbf{\hi{दीवाना}}, \emph{dīvānā}.) An interest? (\textbf{\hi{रुचि}}, \emph{ruci}.) An ambition? (\textbf{\hi{आकांक्षा}}, \emph{ākāṅkṣā}.) A passion? (\textbf{\hi{जुनून}}, \emph{junūn}.) Imagination? (\textbf{\hi{कल्पना}}, \emph{kalpnā}.) A beginning? (\textbf{\hi{शुरुआत}}, \emph{śuruāt}.) An opening? (\textbf{\hi{उद्घाटन}}, \emph{udghāṭan}.) Future? (\textbf{\hi{भावी}}, \emph{bhāvī}.) Life (as one lives it)? (\textbf{\hi{ज़िंदगी}}, \emph{zindagī}.)
\end{tcolorbox}
\section[(dialogue)]{Looking Ahead — once more}

\label{lesson:HI-R559-review-final-559}



\begin{quote}
Say the Hindi for imagination and a beginning.
\end{quote}

\subsection*{Your turn}
\begin{itemize}
\raggedright
  \item \emph{Say it:} imagination, a beginning, an opening, future, life (as one lives it)
\end{itemize}

\begin{tcolorbox}[breakable,colback=teal!4,colframe=teal!35!black,title={Before you move on}]
Imagination? (\textbf{\hi{कल्पना}}, \emph{kalpnā}.) A beginning? (\textbf{\hi{शुरुआत}}, \emph{śuruāt}.) An opening? (\textbf{\hi{उद्घाटन}}, \emph{udghāṭan}.) Future? (\textbf{\hi{भावी}}, \emph{bhāvī}.) Life (as one lives it)? (\textbf{\hi{ज़िंदगी}}, \emph{zindagī}.)
\end{tcolorbox}
